% Generated by tools/scenario-summary.py; do not edit.
\newcommand{\RepairRegressionRows}{%
\bhl{Maint-DefLevel}{整定义重查} & \shortstack{70.9\\{\scriptsize (69.3--74.7)}} & \shortstack{44.0\\{\scriptsize (40.0--53.3)}} & \shortstack{58.0\\{\scriptsize (50.0--60.0)}}\\
\textbf{\bhl{MAVRA}{学习快照回归}} & \shortstack{71.7\\{\scriptsize (69.3--73.3)}} & \shortstack{46.6\\{\scriptsize (40.0--53.8)}} & \shortstack{56.0\\{\scriptsize (50.0--60.0)}}\\
\bhl{MAVRA-CurrData}{当前数据回归} & \shortstack{65.7\\{\scriptsize (58.7--72.0)}} & \shortstack{41.7\textsuperscript{*}\\{\scriptsize (33.3--53.3)}} & \shortstack{60.0\textsuperscript{*}\\{\scriptsize (60.0--60.0)}}\\
\bhl{MAVRA-OracleRef}{真值参照回归} & \shortstack{71.9\\{\scriptsize (68.0--76.0)}} & \shortstack{45.3\\{\scriptsize (40.0--53.3)}} & \shortstack{58.0\\{\scriptsize (50.0--60.0)}}\\
}
\newcommand{\RgMavraStructural}{71.7}
\newcommand{\RgCurrentStructural}{65.7}
\newcommand{\RgOracleStructural}{71.9}

\begin{table}[t]
\caption{\bt{Repair-regression comparison on synthetic retail. Final answer accuracy (\%), including fallback exploration; parentheses give run ranges, not confidence intervals. The five structural scenarios have equal weight.}{合成零售数据的修复回归比较。最终答案正确率（\%），包含回退探索；括号为各次运行的范围，非置信区间。五种结构变化等权。}}
\label{tab:repair-regression}
\centering
\footnotesize
\setlength{\tabcolsep}{3pt}
\renewcommand{\arraystretch}{1.15}
\begin{tabularx}{\linewidth}{@{}Lccc@{}}
\toprule
\bt{Method}{方法} & \bh{Structural}{结构变化} & \bh{Unit}{单位变化} & \bh{Backup}{备份副本}\\
\midrule
\RepairRegressionRows
\bottomrule
\end{tabularx}
\par\smallskip
\begin{minipage}{\linewidth}
\footnotesize
\bt{Five independently learned libraries per method. Asterisks mark four scored runs: provider failures left unit and backup questions unscored in the fifth MAVRA-CurrData run. Pooled values use all scored questions. Repair references are specified in Table~\ref{tab:ladder}.}{每种方法独立学习五个定义库。星号表示四次有计分题的运行：MAVRA-CurrData 第五次运行的单位与备份问题因服务端失败未计分。合并值使用全部计分题。修复参照见表~\ref{tab:ladder}。}
\end{minipage}
\end{table}
